%AttackName: A-GB GCN-1

%Input:
Let $G = (V, E)$ be a graph where $V$ represents the set of nodes and $E$ represents the edges. Let $x \in \mathbb{R}^{|V| \times d}$ be the feature matrix of the nodes, $y$ be the true label associated with the graph, and $f_{\theta}(G)$ be the Graph Convolutional Network (GCN) model. The goal is to generate an adversarial example $G^*$ that misclassifies the input graph.

%Output:
The output of the GB-GCN attack is a modified graph $G^*$ that successfully misclassifies the original input while maintaining structural integrity.

%Formula:
The GB-GCN adversarial attack can be formulated as follows:
1. Initialize the graph:
   $G^{(0)} = G$
2. Set parameters for the optimization process:
   - Define the maximum perturbation size $\epsilon$ and the number of iterations $N$.
3. For each iteration $n = 1$ to $N$:
   - Compute the model's prediction:
   $\hat{y}^{(n)} = f_{\theta}(G^{(n-1)})$
   - Calculate the gradient of the loss function with respect to the node features:
   $g_n = \nabla_x L(f_{\theta}(G^{(n-1)}), y)$
   - Update the node features using the calculated gradient:
   $x^{(n)} = x^{(n-1)} + \delta_n$
   where:
   $\delta_n = -\alpha \cdot \text{sign}(g_n)$
   - Ensure the perturbation stays within the allowed range:
   $x^{(n)} = \text{Clip}_{\mathcal{X}}(x^{(n)})$
   ensuring:
   $\|x^{(n)} - x\|_p \leq \epsilon$

4. The final adversarial graph is:
   $G^* = G^{(N)}$

%Explanation:
Gradient-Based Graph Convolutional Network (GB-GCN) attack leverages gradient information to perturb the node features of a graph in order to mislead a GCN model. The attack iteratively updates the node features based on the gradient of the loss function, ensuring that the perturbations remain within a specified bound. This method effectively generates adversarial examples that can deceive GCNs while preserving the overall structure of the graph, making it a potent technique for evaluating the robustness of graph-based models.